%!TEX root = groundunitbriefing.tex

\chapter{Infantry SAM Units}
\raggedbottom

Infantry SAM units are squads of infantry specialized in the destruction of attacking aircraft with portable SAMs. These units are described below with their counters.

The \hyperlink{table:infantry-sam-unit}{Infantry SAM Unit Table} summarizes the properties of infantry SAM units. In addition to the basic properties shared by all ground units, it gives:
the number of ready, volley, and reload missiles;
the multi-target capability;
whether the unit has quick-reaction capability;
whether the unit has IFF capability;
whether the unit has night IR sights; and
the optical lock-on number and range.
Each squad has one launcher and, unless otherwise specified, one ready missile and two reload missiles. Infantry SAM squads may be attached to infantry, infantry weapons, infantry mortar, or infantry HQ platoons.

%the frequency band, range (in mega-hexes) of the EWR, if present, and whether it has a moving-target indicator (MTI) capability; the frequency band and range (in mega-hexes) of the target-tracking radar, if present; the multi-target capability of the unit; the number of ready missiles and the maximum number of missiles in a volley; whether the unit has a quick-reaction capability; its radar and optical lock-on rolls and (for IR SAMs) its lock-on range in hexes.

The \hyperlink{table:infantry-sam}{Infantry SAM Table} summarizes the properties of the actual missiles:
guidance mode;
launch roll;
turn rate;
flight time;
visibility;
ECCM, chaff, and flare ratings;
whether the missile is active homing (AH) or has a home-on-jam (HOJ) capability;
boost phase in FPs;
base speed and sustainer duration in game turns;
minimum altitude (and modifier for targets at T level);
hit rolls for direct and proximity hits; and
damage ratings for direct and proximity hits.
Most infantry SAMs do not have true proximity fuzes, and so proximity hits in the game correspond to glancing direct hits in reality.

%Some missiles have more than one guidance mode (e.g., the SA-2F has an OG backup mode to its CG main mode).

\subsectionwithcounter{SA-7A Squad}{SA-7A squad}

The 9K32 Strela-2 (SA-7A Grail) was the first Soviet portable SAM. It has a rear-aspect, uncooled IR seeker.
It entered service with Soviet forces in 1968, but quickly was replaced by the improved Strela-2M (SA-7B).%
\note{SA-7A Squad}{
    The launcher and SAM values are from {\AirStr} with significant modifications:
    \begin{itemize}
        \item Zaloga gives an in-service date of 1968.
        \item {\AirStr} gives this a seeker type of I, but that does seem appropriate for an uncooled PbS seeker. Instead, I use a seeker type of E, by analogy with the uncooled seeker in the AIM-9B, and works at about 2 microns.
        \item {\AirStr} gives a turn rate of HT, but the control capability seems to be the same as in the SA-7B and SA-15. I will adopt BT.
        \item {\AirStr} gives this a speed of 8. Zaloga gives a lifetime of 8 seconds and an average speed of 1550 km/h. This corresponds to a range of 5.2 km or about 10 hexes.
    \end{itemize}
}

\subsectionwithcounter{SA-7B Squad}{SA-7B squad}

The 9K32M Strela-2M (SA-7B Grail) is similar to the Strela-2 (SA-7A Grail), but uses a rear-aspect, electrically-cooled IR seeker.
In Warsaw Pact armies, the Strela-2M was sometimes used with the PRP portable RWR receiver that was capable of passively identifying enemy aircraft by their radar emissions.
The Strela-2M entered service with Soviet forces in 1970, replacing the Strela-2.
It has seen combat in almost every conflict since 1970 in which the Soviet Union or its successor states have been participants or supporters.
In particular, the Strela-2M was used heavily by North Vietnamese forces in the 1973 Easter Offensive, by Arab forces in the 1973 Yom Kippur War, by Angolan forces in the South African Border War, by the Afghan Mujahedeen in the Soviet-Afghan War (having been supplied by the CIA), by Iran in the Iran-Iraq War in the 1980s, and by Serbian and Serbian-supported forces in the Yugoslav Wars in the 1990s.%
\note{SA-7B Squad}{
    The launcher and SAM values are from {\AirStr} with significant modifications:
    \begin{itemize}
        \item {\AirStr} gives this a seeker type of I. That seems about right. The seeker was PbS and cooled electrically, like the AIM-9E which also has a type I seeker.
        \item {\AirStr} gives a turn rate of BT, which I adopt.
        \item {\AirStr} gives this a speed of 10. Zaloga states that it has a slightly more powerful rocket motor than the SA-7A, but the mean speed is the same at 1550 km/h. I adopt a speed of 10, like that of the SA-7A.
    \end{itemize}
}

\subsectionwithcounter{SA-14 Squad}{SA-14 squad}

The 9K34 Strela-3 (SA-14 Gremlin) is a development of the earlier Strela-2M.
It has a rear-aspect, cryogen-cooled IR seeker.
In Warsaw Pact armies, the Strela-3 was sometimes used with the PRP portable RWR receiver that was capable of passively identifying enemy aircraft by their radar emissions.
The Strela-3 entered service with the Soviet Army in 1974.
It has seen combat in many conflicts since then in which the Soviet Union or its successor states have been participants or supporters.
In particular, it was used by Iraqi forces in the 1991 Gulf War.%
\note{SA-14 Squad}{
    The launcher and SAM values are from {\AirStr} with significant modifications:
    \begin{itemize}
        \item The seeker is cryogenically cooled at works at about 4 microns (Zaloga). {\AirStr} gives this a type A seeker, but that does seem appropriate; the first type A seekers in Soviet AAMs did not appear until 1985 in the AA-8C. Also, type A seekers have all-aspect capability Zaloga says that the seeker has similar performance to the RIM-43C, but only limited all-aspect capability. A type M seeker seems to be more appropriate.
        \item {\AirStr} gives a turn rate of ET, but the control capability seems to be the same as in the SA-7A and SA-7B. I will adopt BT.
        \item {\AirStr} gives this a speed of 12. Zaloga gives a lifetime of 8 seconds and an average speed of 1440 km/h (as it is slightly heavier than the SA-7). This corresponds to a range of 4.8 km or about 9 hexes.
        \item The warhead is the same as on the SA-7A/7B.
    \end{itemize}
}

\subsectionwithcounter{SA-16 Squad}{SA-16 squad}

The 9K310 Igla-1 (SA-16 Gimlet) uses as single-channel IR seeker similar to that in the SA-14 in a new missile body.
The warhead is similar to the SA-7/14, but also exploded any remaining rocket fuel.
In the game, this is simulated by increasing the damage rating by one if the missile attacks before expending more than half of its FPs (round down).
It incorporates an IFF interrogator, but this is only used in Warsaw Pact armies.
The Igla-1 entered service with the Soviet Army in 1981.
It has been widely exported and have seen combat in many conflicts since then in which the Soviet Union or its successor states have been participants or supporters.
%
\note{SA-16 Squad}{
    The launcher and SAM values are from Malcolm Pipes with significant modifications:
    \begin{itemize}
        \item The warhead of the SA-16/18 is 1.17 kg, identical to the 1.17 kg of the SA-7 and similar to the 1.1 kg of the SA-14. Only later Iglas have heavier warheads. Therefore, I have reduced the damage ratings to match the SA-7.
        \item Carl comments that the warhead explodes any remaining rocket fuel. I increase the damage rating by one if the target is reached in the first half of the FPs.
        \item I have reduced the lock-on range of the SA-16 to match that of the SA-14.
        \item I have increased the flare rating of the SA-16 to match that of the SA-14.
        \item I have given both the SA-14 and SA-16 the same speed and turn rate. Zaloga states that the engagement range increased by about 1 km with respect to the SA-7/14, so I am increasing the speed by 2 to 12. I will give both a turn rate of BT/2, matching Stinger.
    \end{itemize}
}

\subsectionwithcounter{SA-18 Squad}{SA-16 squad}

The 9K38 Igla (SA-18 Grouse) uses a similar missile and launcher similar to the 9K310 Igla-1 (SA-16), but with a new a dual-channel IR seeker to reduce its susceptibility to flares.
Again, the warhead is similar to the SA-7/14, but also exploded any remaining rocket fuel.
In the game, this is simulated by increasing the damage rating by one if the missile attacks before expending more than half of its FPs (round down).
It incorporates an IFF interrogator, but this is only used in Warsaw Pact armies.
The Igla entered service with the Soviet Army in 1983.
It has been widely exported and have seen combat in many conflicts since then in which the Soviet Union or its successor states have been participants or supporters.
It was used by Iraqi forces in the 1991 Gulf War, by the Peruvian Army in the 1995 Cenepa War, and by Serbian forces in the 1999 Kosovo War.%
\note{SA-18 Squad}{
    The launcher and SAM values are from Malcolm Pipes with some modifications:
    \begin{itemize}
        \item The warhead of the SA-16/18 is 1.17 kg, identical to the 1.17 kg of the SA-7 and similar to the 1.1 kg of the SA-14. Only later Iglas had heavier warheads. Therefore, I have reduced the damage ratings to match the SA-7.
        \item Carl comments that the warhead explodes any remaining rocket fuel. I increase the damage rating by one if the target is reached in the first half of the FPs.
        \item I have given both the SA-14 and SA-16 the same speed and turn rate. Zaloga states that the engagement range increased by about 1 km with respect to the SA-7/14, so I am increasing the speed by 2 to 12. I will give both a turn rate of BT/2, matching Stinger.
    \end{itemize}
}

\subsectionwithcounter{FIM-43C Squad}{FIM-43C squad}

The FIM-43C Redeye Block III was the first US infantry SAM to reach production and enter service.
It has a rear-aspect, cryogen-cooled IR seeker.
The FIM-43C entered service with the US Army in 1968 and subsequently was used by Australia, Chad, Denmark, West Germany, Greece, Israel, Jordan, Saudi Arabia, Somalia, Sudan, Sweden, and Turkey. It was also supplied by the CIA to the Afghan Mujahedeen and the Nicaraguan Contras.%
\note{FIM-43C Squad}{
    The launcher and SAM values are from {\AirStr}.
}

\subsectionwithfourcounters{FIM-92A/B/C/D Squad}{FIM-92A squad}{FIM-92B squad}{FIM-92C squad}{FIM-92D squad}

The FIM-92A Stinger has an all-angle, second-generation, cryogen-cooled IR seeker.
Its warhead weight 3 kg and was significantly more lethal than in earlier missiles.
The FIM-92B Stinger POST adds a UV seeker to distinguish flares from aircraft.
The FIM-92C/D Stinger RMP have software updates to further improve the performance against flares.
All have an IFF interrogator to help avoid fratricide.
The FIM-92A entered service with the US Army in 1981, the FIM-92B in 1986, the FIM-92C in 1989, and the FIM-92D in 1992. They were exported to dozens of US allies including the Afghan Mujahedeen and the Angolan UNITA. They saw service in the 1982 South Atlantic War, the Soviet-Afghan War, the Angolan Civil War, the Libyan Invasion of Chad, the Tajik Civil War, the Chechen War, the Sri Lankan Civil War, the Syrian Civil War, and the Russian Invasion of Ukraine.%
\note{FIM-92A/B/C/D Squad}{
    The launcher and SAM values for the A model are from {\AirStr}. Those for the B/C/D models are from Malcolm Pipes compendium.
    \begin{itemize}
        \item I have increased the flare susceptibility of the FIM-92A from 2 to 3. It did not seem to have significant protection.
        \item Wikipedia gives an effective range of 5 miles, which is equivalent to a speed of 15. I have left the speed at 14, as that is close enough.
        \item Wikipedia gives a peak speed of 1930 mp and a self-destruct time of 17 seconds.
        \item {\AirStr} has the visibility as 2, but looking at photos I don't see much differ
    \end{itemize}
}

\subsectionwithcounter{Blowpipe Squad}{Blowpipe squad}

Blowpipe is an MCLOS optically-guided missile.
While in theory this gave it an all-aspect capability, in practice guiding the missile to the target in combat proved to be very difficult.
Blowpipe entered service with the British Army and Royal Marines in 1975.
It was also used by the Argentine Army, Canadian Army, Chilean Army and Navy, Ecuadorean Army, Guatemalan Army, Portuguese Army, Malawian Army, Malaysian Army, Nigerian Army, Omani Army, Royal Thai Air Force and Army, and the UAE Army.
It saw combat in the 1982 South Atlantic War with both sides and had a low success rate.
In 1986, it was trialed by the Afghan Mujahedeen, but again was not considered to be effective.
Finally, it was also used in the 1995 Cenepa War by Ecuador.%
\note{Blowpipe Squad}{
    The launcher and SAM values are from {\AirStr}.
}

\subsectionwithcounter{Javelin Squad}{Javelin squad}

Given the poor performance of Blowpipe in the 1982 South Atlantic War, Javelin was quickly developed as an upgraded version of Blowpipe with SACLOS guidance in place of MCLOS guidance.
Javelin replaced Blowpipe in the British Army and Royal Marines in 1984.
It was also used by the Botswana Army, the Canadian Army, the Malaysian Army, the Peruvian Army, and the South Korean Army.%
\note{Javelin Squad}{
    The launcher and SAM values for Javelin are from {\AirStr}.

    Javelin is also called “Javelin GL”.

    I have seen mention of a Javelin LML; did anyone use it or was it just a prototype?
}

\subsectionwithcounter{Starburst Squad}{Starburst squad}

Starburst is Javelin with the radio-link replaced by a with laser link.
Starburst replaced Javelin in the British Army and Royal Marines in 1989.
It was also used by the Canadian Army, the Malaysian Army, the Qatari Army, and the Royal Thai Army.
Starburst served with the British Army in the 1991 Gulf War, but was not used in combat.%
\note{Starburst Squad}{
    The launcher and SAM values for Starburst are just those for Javelin but with the OG guidance replaced by LG guidance.

    Starburst was originally called “Javelin S15”.
}

\subsectionwithcounter{Starstreak Squad}{Starstreak squad}

Starstreak is a high-speed laser-guided missile. As the missile approaches the target, it launches three independent explosive darts. The darts have contact fuzes, but not proximity fuzes.
Starstreak replaced Javelin in the British Army and Royal Marines in 2000 and is also used by the South African Army and the Armed Forces of Ukraine.
Starstreak served with the British Army in the 2003 Invasion of Iraq, but was not used in combat. It has seen combat in the Russian Invasion of Ukraine.%
\note{Starstreak Squad}{
    The launcher and SAM values are from Malcolm Pipes. The real darts do not have proximity fuzes, but the missile in the game can obtain a proximity hit. Perhaps a proximity hit here means a hit with only one dart?
}

\subsectionwithcounter{Starstreak LML Squad}{Starstreak LML squad}

The Starstreak Light Multiple Launcher (LML) has three ready-to-fire Starstreak missiles on a pedestal mount with a single sighting unit. The launcher and its associated equipment are heavy and will normally be transported by a vehicle unit.
Starstreak LML entered service with the British Army and Royal Marines in 2000.
It is also used by Indonesia, Malaysia, and South Africa.%
\note{Starstreak LML Squad}{
    The SAM values are from Malcolm Pipes. I have guessed the launcher values are the same as for shoulder-launched version, just that the launcher has three ready missiles. I don't know if the launcher has reloads. Is the launcher really portable or does it have to be transported by a vehicle?
}

\subsectionwithcounter{Mistral Squad}{Mistral squad}

The Mistral is a French IR-homing missile on a pedestal mount. The missile has a considerably larger warhead than either the Stinger or Igla missiles, but this comes at a cost in portability. The launcher and its associated equipment are heavy and will normally be transported by a vehicle unit.
The Mistral entered service with the French Army in 1990. It was subsequently exported to many other countries.%
\note{Mistral Squad}{
    The launcher and SAM values are from {\AirStr}. Is the launcher really portable or does it have to be transported by a vehicle?
}

\subsectionwithcounter{RBS 70 Squad}{RBS 70 squad}
The RBS 70 is a Swedish laser-guided missile on a pedestal mount. The launcher and its associated equipment are heavy and will normally be transported by a vehicle unit.
The RBS 70 entered service with the Swedish Army in 1977. It was subsequently used by Argentina, Australia, Bahrain, Brazil, the Czech Republic, Iran, Ireland, Lithuania, Norway, Pakistan, Singapore, Thailand, Tunisia, the UAE, Ukraine, and Venezuela.%
\note{RBS 70 Squad}{
    For the RBS 70: The launcher and SAM values are from {\AirStr}. Is the launcher really portable or does it have to be transported by a vehicle?
}

\subsectionwithcounter{Ayn-al-Saqr Squad}{Ayn-al-Saqr squad}

The Ayn-al-Saqr (“Eye of the Hawk”) missile is a  licensed version of the 9M32M Strela-2m (SA-7B Grail) manufactured in Egypt. It equips one version of the Sinai 23 air-defense vehicle.
It entered service with the Egyptian Army in 1984.%
\note{Ayn-al-Saqr Squad}{
    The launcher and SAM values are taken to be equal to those of the SA-7B.
}

\begin{table*}[p]
    \tablecaption{table:infantry-sam-unit}{Infantry SAM Unit Table}
    \centering
    \footnotesize
    \input{table-infantry-sam-units.tex}
\end{table*}

\begin{table*}[p]
    \tablecaption{table:infantry-sam}{Infantry SAM Table}
    \centering
    \footnotesize
    \input{table-infantry-sam.tex}
\end{table*}
